
\documentclass{article}
\usepackage{colortbl}
\usepackage{makecell}
\usepackage{multirow}
\usepackage{supertabular}

\begin{document}

\newcounter{utterance}

\centering \large Interaction Transcript for game `bbh', experiment `navigate', episode 71 with reflexion{-}llama{-}3.1{-}8b{-}instruct.
\vspace{24pt}

{ \footnotesize  \setcounter{utterance}{1}
\setlength{\tabcolsep}{0pt}
\begin{supertabular}{c@{$\;$}|p{.15\linewidth}@{}p{.15\linewidth}p{.15\linewidth}p{.15\linewidth}p{.15\linewidth}p{.15\linewidth}}
    \# & \multicolumn{2}{c}{Player} && \multicolumn{2}{c}{Game Master} \\
    \hline

    \theutterance \stepcounter{utterance}  
    & & & \multicolumn{4}{p{0.6\linewidth}}{
        \cellcolor[rgb]{0.9,0.9,0.9}{
            \makecell[{{p{\linewidth}}}]{
                \texttt{\tiny{[P1$\langle$GM]}}
                \texttt{You have attempted this same task before and did not succeed. These are the plans you wrote after those attempts.} \\
\texttt{Trial \#1: New plan: } \\
\\ 
\texttt{In my previous attempt, I correctly analyzed the instructions but failed to consider the implications of the "always face forward" instruction. I should have taken into account that the "backward" steps would be to the left or right, not directly opposite the direction I'm facing. However, I also failed to recognize that the 4 steps backward would effectively cancel out the 4 steps forward that would have been taken if I were facing the opposite direction. } \\
\\ 
\texttt{To improve my strategy, I will pay closer attention to the instructions and consider the implications of each action. In this case, I should have recognized that the 4 steps backward would effectively cancel out the 4 steps forward, resulting in a return to the starting point. } \\
\\ 
\texttt{In general, I will strive to be more thorough in my analysis and consider all possible implications of the instructions. I will also look for patterns and relationships between the actions that can help me determine the outcome.} \\
\\ 
\\ 
\texttt{Given a series of navigation instructions, determine whether one would end up back at the starting point.} \\
\\ 
\texttt{Q: If you follow these instructions, do you return to the starting point? Turn left. Turn around. Turn left. Take 7 steps. Take 2 steps. Take 4 steps. Take 8 steps.} \\
\texttt{Options:} \\
\texttt{{-} Yes} \\
\texttt{{-} No} \\
\texttt{A: No} \\
\\ 
\texttt{Q: If you follow these instructions, do you return to the starting point? Turn around. Take 1 step. Take 6 steps. Turn around. Take 6 steps. Take 9 steps. Take 1 step.} \\
\texttt{Options:} \\
\texttt{{-} Yes} \\
\texttt{{-} No} \\
\texttt{A: No} \\
\\ 
\texttt{Q: If you follow these instructions, do you return to the starting point? Always face forward. Take 2 steps right. Take 9 steps left. Take 7 steps right.} \\
\texttt{Options:} \\
\texttt{{-} Yes} \\
\texttt{{-} No} \\
\texttt{A: Yes} \\
\\ 
\texttt{Q: If you follow these instructions, do you return to the starting point? Always face forward. Take 6 steps forward. Take 2 steps backward. Take 4 steps backward.} \\
\texttt{Options:} \\
\texttt{{-} Yes} \\
\texttt{{-} No} \\
\texttt{A:} \\
            }
        }
    }
     \\ \\

    \theutterance \stepcounter{utterance}  
    & \multicolumn{4}{p{0.6\linewidth}}{
        \cellcolor[rgb]{0.561,0.737,0.561}{
            \makecell[{{p{\linewidth}}}]{
                \texttt{\tiny{[P1$\rangle$GM]}}
                \texttt{Based on the given instructions and the analysis from my previous attempt, I will apply the same strategy to determine whether one would end up back at the starting point.} \\
\\ 
\texttt{Q: If you follow these instructions, do you return to the starting point? Always face forward. Take 6 steps forward. Take 2 steps backward. Take 4 steps backward.} \\
\\ 
\texttt{Since the instructions state "Always face forward," the "backward" steps would be to the left or right, not directly opposite the direction I'm facing. } \\
\\ 
\texttt{When taking 2 steps backward, it would be 2 steps to the left or right. } \\
\texttt{When taking 4 steps backward, it would be 4 steps to the left or right. } \\
\\ 
\texttt{Considering the 2 steps backward and 4 steps backward, the total movement to the left or right would be 6 steps. } \\
\\ 
\texttt{Since the initial movement was 6 steps forward, the net movement would be 0 steps forward or backward. } \\
\\ 
\texttt{However, the 6 steps forward would be canceled out by the 6 steps to the left or right. } \\
\\ 
\texttt{Therefore, the net movement would be 6 steps to the left or right, not back to the starting point.} \\
\\ 
\texttt{Answer: No} \\
            }
        }
    }
    & & \\ \\

    \theutterance \stepcounter{utterance}  
    & & & \multicolumn{2}{p{0.3\linewidth}}{
        \cellcolor[rgb]{0.9,0.9,0.9}{
            \makecell[{{p{\linewidth}}}]{
                \texttt{\tiny{[GM$|$GM]}}
                \texttt{Target: Yes} \\
            }
        }
    }
    & & \\ \\

    \theutterance \stepcounter{utterance}  
    & & & \multicolumn{2}{p{0.3\linewidth}}{
        \cellcolor[rgb]{0.9,0.9,0.9}{
            \makecell[{{p{\linewidth}}}]{
                \texttt{\tiny{[GM$|$GM]}}
                \texttt{game\_result = LOSE} \\
            }
        }
    }
    & & \\ \\

    \theutterance \stepcounter{utterance}  
    & & \multicolumn{4}{p{0.6\linewidth}}{
        \cellcolor[rgb]{0.95,0.95,0.95}{
            \makecell[{{p{\linewidth}}}]{
                \texttt{\tiny{[P1$\leftrightarrow$P]}}
                \texttt{Post{-}episode Reflexion (not sent to GM):} \\
\texttt{New plan: } \\
\\ 
\texttt{In my previous attempt, I correctly analyzed the instructions but failed to recognize that the 4 steps backward would effectively cancel out the 4 steps forward, resulting in a return to the starting point. However, I also failed to consider the fact that the 2 steps backward would also cancel out the 2 steps forward that would have been taken if I were facing the opposite direction. } \\
\\ 
\texttt{To improve my strategy, I will pay closer attention to the instructions and consider the implications of each action. In this case, I should have recognized that the 2 steps backward and 4 steps backward would effectively cancel out the 2 steps forward and 4 steps forward, resulting in a return to the starting point. } \\
\\ 
\texttt{In general, I will strive to be more thorough in my analysis and consider all possible implications of the instructions. I will also look for patterns and relationships between the actions that can help me determine the outcome.} \\
            }
        }
    }
    & \\ \\

\end{supertabular}
}

\end{document}
